%% itbdisertasi-layout.tex
%% ---------------------------------------------------------------------------
%% Single source of truth for ITB dissertation layout tokens.
%%
%% Loaded by |itbdisertasi.cls| via |%% itbdisertasi-layout.tex
%% ---------------------------------------------------------------------------
%% Single source of truth for ITB dissertation layout tokens.
%%
%% Loaded by |itbdisertasi.cls| via |%% itbdisertasi-layout.tex
%% ---------------------------------------------------------------------------
%% Single source of truth for ITB dissertation layout tokens.
%%
%% Loaded by |itbdisertasi.cls| via |%% itbdisertasi-layout.tex
%% ---------------------------------------------------------------------------
%% Single source of truth for ITB dissertation layout tokens.
%%
%% Loaded by |itbdisertasi.cls| via |\input{itbdisertasi-layout.tex}|.
%% Compile from the directory that contains this file (same folder as
%% |disertasi.tex|) so the input resolves.
%%
%% Sections:
%%   (1) Page geometry — mirror margins (Pedoman / Juknis SPs ITB)
%%   (2) Heading vertical rhythm — titlesec (chapter aligns to 3 cm rule)
%%   (3) TOC / LOF / LOT — tocloft skips and list title vertical tweak
%%   (4) Sampul & halaman pengesahan — vertical gaps (juknis diagram)
%%
%% To adjust globally: change the macro or skip value here only, then
%% rebuild. For empirical fine-tuning of chapter vs.\ topskip offset,
%% edit |\ITBLayoutChapterTitleAbove| (negative = pull title up).
%% ---------------------------------------------------------------------------

% ---------------------------------------------------------------------------
% (1) Page geometry — applied via |\ITBApplyPageGeometry| after
%     |\RequirePackage{geometry}| in the class file.
% ---------------------------------------------------------------------------
\newcommand{\ITBGeomPaper}{a4paper}
\newcommand{\ITBGeomInner}{4cm}
\newcommand{\ITBGeomOuter}{3cm}
\newcommand{\ITBGeomTop}{3cm}
\newcommand{\ITBGeomBottom}{3cm}
\newcommand{\ITBGeomHeadHeight}{15pt}
\newcommand{\ITBGeomHeadSep}{12pt}
\newcommand{\ITBGeomFootSkip}{30pt}

\newcommand{\ITBApplyPageGeometry}{%
  \geometry{%
    \ITBGeomPaper,%
    inner=\ITBGeomInner,%
    outer=\ITBGeomOuter,%
    top=\ITBGeomTop,%
    bottom=\ITBGeomBottom,%
    headheight=\ITBGeomHeadHeight,%
    headsep=\ITBGeomHeadSep,%
    footskip=\ITBGeomFootSkip,%
    twoside%
  }%
}

% ---------------------------------------------------------------------------
% (2) Chapter / section vertical spacing (titlesec + report baseline).
%     Chapter: single line "Bab I  Judul" per juknis Mei 2019 hal.\ 6.
% ---------------------------------------------------------------------------
\newskip\ITBLayoutChapterTitleAbove
\newskip\ITBLayoutChapterTitleBelow
\newskip\ITBLayoutSectionTitleAbove
\newskip\ITBLayoutSectionTitleBelow
\newskip\ITBLayoutSubsectionTitleAbove
\newskip\ITBLayoutSubsectionTitleBelow
\newskip\ITBLayoutSubsubsectionTitleAbove
\newskip\ITBLayoutSubsubsectionTitleBelow

\ITBLayoutChapterTitleAbove=-36pt plus 0pt minus 0pt
\ITBLayoutChapterTitleBelow=18pt plus 0pt minus 0pt
\ITBLayoutSectionTitleAbove=18pt plus 0pt minus 0pt
\ITBLayoutSectionTitleBelow=6pt plus 0pt minus 0pt
\ITBLayoutSubsectionTitleAbove=14pt plus 0pt minus 0pt
\ITBLayoutSubsectionTitleBelow=4pt plus 0pt minus 0pt
\ITBLayoutSubsubsectionTitleAbove=12pt plus 0pt minus 0pt
\ITBLayoutSubsubsectionTitleBelow=4pt plus 0pt minus 0pt

% ---------------------------------------------------------------------------
% (3) Daftar isi / gambar / tabel — tocloft vertical density + title skip
% ---------------------------------------------------------------------------
\newcommand{\ITBLayoutCftBeforeChapSkip}{0.2em}
\newcommand{\ITBLayoutCftBeforeSecSkip}{0.1em}
\newcommand{\ITBLayoutCftBeforeSubsecSkip}{0.1em}
\newcommand{\ITBLayoutCftBeforeFigSkip}{0.1em}
\newcommand{\ITBLayoutCftBeforeTabSkip}{0.1em}
\newlength{\ITBLayoutCftBeforeTocTitle}
\newlength{\ITBLayoutCftBeforeLoFTitle}
\newlength{\ITBLayoutCftBeforeLoTTitle}
\newlength{\ITBLayoutCftAfterTocTitle}
\newlength{\ITBLayoutCftAfterLoFTitle}
\newlength{\ITBLayoutCftAfterLoTTitle}

\setlength{\ITBLayoutCftBeforeTocTitle}{-20pt}
\setlength{\ITBLayoutCftBeforeLoFTitle}{-20pt}
\setlength{\ITBLayoutCftBeforeLoTTitle}{-20pt}
\setlength{\ITBLayoutCftAfterTocTitle}{18pt}
\setlength{\ITBLayoutCftAfterLoFTitle}{18pt}
\setlength{\ITBLayoutCftAfterLoTTitle}{18pt}

\newcommand{\ITBApplyTocloftLayout}{%
  \setlength{\cftbeforechapskip}{\ITBLayoutCftBeforeChapSkip}%
  \setlength{\cftbeforesecskip}{\ITBLayoutCftBeforeSecSkip}%
  \setlength{\cftbeforesubsecskip}{\ITBLayoutCftBeforeSubsecSkip}%
  \setlength{\cftbeforefigskip}{\ITBLayoutCftBeforeFigSkip}%
  \setlength{\cftbeforetabskip}{\ITBLayoutCftBeforeTabSkip}%
  \setlength{\cftbeforetoctitleskip}{\ITBLayoutCftBeforeTocTitle}%
  \setlength{\cftbeforeloftitleskip}{\ITBLayoutCftBeforeLoFTitle}%
  \setlength{\cftbeforelottitleskip}{\ITBLayoutCftBeforeLoTTitle}%
  \setlength{\cftaftertoctitleskip}{\ITBLayoutCftAfterTocTitle}%
  \setlength{\cftafterloftitleskip}{\ITBLayoutCftAfterLoFTitle}%
  \setlength{\cftafterlottitleskip}{\ITBLayoutCftAfterLoTTitle}%
}

% ---------------------------------------------------------------------------
% (4) Sampul & halaman pengesahan — vertical gaps (juknis diagram)
% ---------------------------------------------------------------------------
\newcommand{\ITBSampulTopGap}{0.2cm}
\newcommand{\ITBSampulAfterTitle}{1.0cm}
\newcommand{\ITBSampulAfterDisertasi}{0.5cm}
\newcommand{\ITBSampulAfterKaryaTulis}{0.9cm}
\newcommand{\ITBSampulAfterOleh}{0.3cm}
\newcommand{\ITBSampulAfterNimBlock}{0.2cm}
\newcommand{\ITBSampulAfterProdi}{0.9cm}
\newcommand{\ITBSampulLogoHeight}{3.5cm}
\newcommand{\ITBSampulAfterInstitusi}{0.15cm}

\newcommand{\ITBPengesahanAfterHeading}{0.8cm}
\newcommand{\ITBPengesahanAfterTitle}{0.8cm}
\newcommand{\ITBPengesahanAfterAuthorBlock}{0.9cm}
\newcommand{\ITBPengesahanAfterMenyetujui}{0.3cm}
\newcommand{\ITBPengesahanAfterTanggal}{0.6cm}
\newcommand{\ITBPengesahanAfterKetuaLabel}{1.5cm}
\newcommand{\ITBPengesahanAfterKetuaRule}{0.8cm}
\newcommand{\ITBPengesahanAnggotaRuleDrop}{1.5cm}

\newcommand{\ITBPedomanBetweenParagraphs}{0.5cm}
|.
%% Compile from the directory that contains this file (same folder as
%% |disertasi.tex|) so the input resolves.
%%
%% Sections:
%%   (1) Page geometry — mirror margins (Pedoman / Juknis SPs ITB)
%%   (2) Heading vertical rhythm — titlesec (chapter aligns to 3 cm rule)
%%   (3) TOC / LOF / LOT — tocloft skips and list title vertical tweak
%%   (4) Sampul & halaman pengesahan — vertical gaps (juknis diagram)
%%
%% To adjust globally: change the macro or skip value here only, then
%% rebuild. For empirical fine-tuning of chapter vs.\ topskip offset,
%% edit |\ITBLayoutChapterTitleAbove| (negative = pull title up).
%% ---------------------------------------------------------------------------

% ---------------------------------------------------------------------------
% (1) Page geometry — applied via |\ITBApplyPageGeometry| after
%     |\RequirePackage{geometry}| in the class file.
% ---------------------------------------------------------------------------
\newcommand{\ITBGeomPaper}{a4paper}
\newcommand{\ITBGeomInner}{4cm}
\newcommand{\ITBGeomOuter}{3cm}
\newcommand{\ITBGeomTop}{3cm}
\newcommand{\ITBGeomBottom}{3cm}
\newcommand{\ITBGeomHeadHeight}{15pt}
\newcommand{\ITBGeomHeadSep}{12pt}
\newcommand{\ITBGeomFootSkip}{30pt}

\newcommand{\ITBApplyPageGeometry}{%
  \geometry{%
    \ITBGeomPaper,%
    inner=\ITBGeomInner,%
    outer=\ITBGeomOuter,%
    top=\ITBGeomTop,%
    bottom=\ITBGeomBottom,%
    headheight=\ITBGeomHeadHeight,%
    headsep=\ITBGeomHeadSep,%
    footskip=\ITBGeomFootSkip,%
    twoside%
  }%
}

% ---------------------------------------------------------------------------
% (2) Chapter / section vertical spacing (titlesec + report baseline).
%     Chapter: single line "Bab I  Judul" per juknis Mei 2019 hal.\ 6.
% ---------------------------------------------------------------------------
\newskip\ITBLayoutChapterTitleAbove
\newskip\ITBLayoutChapterTitleBelow
\newskip\ITBLayoutSectionTitleAbove
\newskip\ITBLayoutSectionTitleBelow
\newskip\ITBLayoutSubsectionTitleAbove
\newskip\ITBLayoutSubsectionTitleBelow
\newskip\ITBLayoutSubsubsectionTitleAbove
\newskip\ITBLayoutSubsubsectionTitleBelow

\ITBLayoutChapterTitleAbove=-36pt plus 0pt minus 0pt
\ITBLayoutChapterTitleBelow=18pt plus 0pt minus 0pt
\ITBLayoutSectionTitleAbove=18pt plus 0pt minus 0pt
\ITBLayoutSectionTitleBelow=6pt plus 0pt minus 0pt
\ITBLayoutSubsectionTitleAbove=14pt plus 0pt minus 0pt
\ITBLayoutSubsectionTitleBelow=4pt plus 0pt minus 0pt
\ITBLayoutSubsubsectionTitleAbove=12pt plus 0pt minus 0pt
\ITBLayoutSubsubsectionTitleBelow=4pt plus 0pt minus 0pt

% ---------------------------------------------------------------------------
% (3) Daftar isi / gambar / tabel — tocloft vertical density + title skip
% ---------------------------------------------------------------------------
\newcommand{\ITBLayoutCftBeforeChapSkip}{0.2em}
\newcommand{\ITBLayoutCftBeforeSecSkip}{0.1em}
\newcommand{\ITBLayoutCftBeforeSubsecSkip}{0.1em}
\newcommand{\ITBLayoutCftBeforeFigSkip}{0.1em}
\newcommand{\ITBLayoutCftBeforeTabSkip}{0.1em}
\newlength{\ITBLayoutCftBeforeTocTitle}
\newlength{\ITBLayoutCftBeforeLoFTitle}
\newlength{\ITBLayoutCftBeforeLoTTitle}
\newlength{\ITBLayoutCftAfterTocTitle}
\newlength{\ITBLayoutCftAfterLoFTitle}
\newlength{\ITBLayoutCftAfterLoTTitle}

\setlength{\ITBLayoutCftBeforeTocTitle}{-20pt}
\setlength{\ITBLayoutCftBeforeLoFTitle}{-20pt}
\setlength{\ITBLayoutCftBeforeLoTTitle}{-20pt}
\setlength{\ITBLayoutCftAfterTocTitle}{18pt}
\setlength{\ITBLayoutCftAfterLoFTitle}{18pt}
\setlength{\ITBLayoutCftAfterLoTTitle}{18pt}

\newcommand{\ITBApplyTocloftLayout}{%
  \setlength{\cftbeforechapskip}{\ITBLayoutCftBeforeChapSkip}%
  \setlength{\cftbeforesecskip}{\ITBLayoutCftBeforeSecSkip}%
  \setlength{\cftbeforesubsecskip}{\ITBLayoutCftBeforeSubsecSkip}%
  \setlength{\cftbeforefigskip}{\ITBLayoutCftBeforeFigSkip}%
  \setlength{\cftbeforetabskip}{\ITBLayoutCftBeforeTabSkip}%
  \setlength{\cftbeforetoctitleskip}{\ITBLayoutCftBeforeTocTitle}%
  \setlength{\cftbeforeloftitleskip}{\ITBLayoutCftBeforeLoFTitle}%
  \setlength{\cftbeforelottitleskip}{\ITBLayoutCftBeforeLoTTitle}%
  \setlength{\cftaftertoctitleskip}{\ITBLayoutCftAfterTocTitle}%
  \setlength{\cftafterloftitleskip}{\ITBLayoutCftAfterLoFTitle}%
  \setlength{\cftafterlottitleskip}{\ITBLayoutCftAfterLoTTitle}%
}

% ---------------------------------------------------------------------------
% (4) Sampul & halaman pengesahan — vertical gaps (juknis diagram)
% ---------------------------------------------------------------------------
\newcommand{\ITBSampulTopGap}{0.2cm}
\newcommand{\ITBSampulAfterTitle}{1.0cm}
\newcommand{\ITBSampulAfterDisertasi}{0.5cm}
\newcommand{\ITBSampulAfterKaryaTulis}{0.9cm}
\newcommand{\ITBSampulAfterOleh}{0.3cm}
\newcommand{\ITBSampulAfterNimBlock}{0.2cm}
\newcommand{\ITBSampulAfterProdi}{0.9cm}
\newcommand{\ITBSampulLogoHeight}{3.5cm}
\newcommand{\ITBSampulAfterInstitusi}{0.15cm}

\newcommand{\ITBPengesahanAfterHeading}{0.8cm}
\newcommand{\ITBPengesahanAfterTitle}{0.8cm}
\newcommand{\ITBPengesahanAfterAuthorBlock}{0.9cm}
\newcommand{\ITBPengesahanAfterMenyetujui}{0.3cm}
\newcommand{\ITBPengesahanAfterTanggal}{0.6cm}
\newcommand{\ITBPengesahanAfterKetuaLabel}{1.5cm}
\newcommand{\ITBPengesahanAfterKetuaRule}{0.8cm}
\newcommand{\ITBPengesahanAnggotaRuleDrop}{1.5cm}

\newcommand{\ITBPedomanBetweenParagraphs}{0.5cm}
|.
%% Compile from the directory that contains this file (same folder as
%% |disertasi.tex|) so the input resolves.
%%
%% Sections:
%%   (1) Page geometry — mirror margins (Pedoman / Juknis SPs ITB)
%%   (2) Heading vertical rhythm — titlesec (chapter aligns to 3 cm rule)
%%   (3) TOC / LOF / LOT — tocloft skips and list title vertical tweak
%%   (4) Sampul & halaman pengesahan — vertical gaps (juknis diagram)
%%
%% To adjust globally: change the macro or skip value here only, then
%% rebuild. For empirical fine-tuning of chapter vs.\ topskip offset,
%% edit |\ITBLayoutChapterTitleAbove| (negative = pull title up).
%% ---------------------------------------------------------------------------

% ---------------------------------------------------------------------------
% (1) Page geometry — applied via |\ITBApplyPageGeometry| after
%     |\RequirePackage{geometry}| in the class file.
% ---------------------------------------------------------------------------
\newcommand{\ITBGeomPaper}{a4paper}
\newcommand{\ITBGeomInner}{4cm}
\newcommand{\ITBGeomOuter}{3cm}
\newcommand{\ITBGeomTop}{3cm}
\newcommand{\ITBGeomBottom}{3cm}
\newcommand{\ITBGeomHeadHeight}{15pt}
\newcommand{\ITBGeomHeadSep}{12pt}
\newcommand{\ITBGeomFootSkip}{30pt}

\newcommand{\ITBApplyPageGeometry}{%
  \geometry{%
    \ITBGeomPaper,%
    inner=\ITBGeomInner,%
    outer=\ITBGeomOuter,%
    top=\ITBGeomTop,%
    bottom=\ITBGeomBottom,%
    headheight=\ITBGeomHeadHeight,%
    headsep=\ITBGeomHeadSep,%
    footskip=\ITBGeomFootSkip,%
    twoside%
  }%
}

% ---------------------------------------------------------------------------
% (2) Chapter / section vertical spacing (titlesec + report baseline).
%     Chapter: single line "Bab I  Judul" per juknis Mei 2019 hal.\ 6.
% ---------------------------------------------------------------------------
\newskip\ITBLayoutChapterTitleAbove
\newskip\ITBLayoutChapterTitleBelow
\newskip\ITBLayoutSectionTitleAbove
\newskip\ITBLayoutSectionTitleBelow
\newskip\ITBLayoutSubsectionTitleAbove
\newskip\ITBLayoutSubsectionTitleBelow
\newskip\ITBLayoutSubsubsectionTitleAbove
\newskip\ITBLayoutSubsubsectionTitleBelow

\ITBLayoutChapterTitleAbove=-36pt plus 0pt minus 0pt
\ITBLayoutChapterTitleBelow=18pt plus 0pt minus 0pt
\ITBLayoutSectionTitleAbove=18pt plus 0pt minus 0pt
\ITBLayoutSectionTitleBelow=6pt plus 0pt minus 0pt
\ITBLayoutSubsectionTitleAbove=14pt plus 0pt minus 0pt
\ITBLayoutSubsectionTitleBelow=4pt plus 0pt minus 0pt
\ITBLayoutSubsubsectionTitleAbove=12pt plus 0pt minus 0pt
\ITBLayoutSubsubsectionTitleBelow=4pt plus 0pt minus 0pt

% ---------------------------------------------------------------------------
% (3) Daftar isi / gambar / tabel — tocloft vertical density + title skip
% ---------------------------------------------------------------------------
\newcommand{\ITBLayoutCftBeforeChapSkip}{0.2em}
\newcommand{\ITBLayoutCftBeforeSecSkip}{0.1em}
\newcommand{\ITBLayoutCftBeforeSubsecSkip}{0.1em}
\newcommand{\ITBLayoutCftBeforeFigSkip}{0.1em}
\newcommand{\ITBLayoutCftBeforeTabSkip}{0.1em}
\newlength{\ITBLayoutCftBeforeTocTitle}
\newlength{\ITBLayoutCftBeforeLoFTitle}
\newlength{\ITBLayoutCftBeforeLoTTitle}
\newlength{\ITBLayoutCftAfterTocTitle}
\newlength{\ITBLayoutCftAfterLoFTitle}
\newlength{\ITBLayoutCftAfterLoTTitle}

\setlength{\ITBLayoutCftBeforeTocTitle}{-20pt}
\setlength{\ITBLayoutCftBeforeLoFTitle}{-20pt}
\setlength{\ITBLayoutCftBeforeLoTTitle}{-20pt}
\setlength{\ITBLayoutCftAfterTocTitle}{18pt}
\setlength{\ITBLayoutCftAfterLoFTitle}{18pt}
\setlength{\ITBLayoutCftAfterLoTTitle}{18pt}

\newcommand{\ITBApplyTocloftLayout}{%
  \setlength{\cftbeforechapskip}{\ITBLayoutCftBeforeChapSkip}%
  \setlength{\cftbeforesecskip}{\ITBLayoutCftBeforeSecSkip}%
  \setlength{\cftbeforesubsecskip}{\ITBLayoutCftBeforeSubsecSkip}%
  \setlength{\cftbeforefigskip}{\ITBLayoutCftBeforeFigSkip}%
  \setlength{\cftbeforetabskip}{\ITBLayoutCftBeforeTabSkip}%
  \setlength{\cftbeforetoctitleskip}{\ITBLayoutCftBeforeTocTitle}%
  \setlength{\cftbeforeloftitleskip}{\ITBLayoutCftBeforeLoFTitle}%
  \setlength{\cftbeforelottitleskip}{\ITBLayoutCftBeforeLoTTitle}%
  \setlength{\cftaftertoctitleskip}{\ITBLayoutCftAfterTocTitle}%
  \setlength{\cftafterloftitleskip}{\ITBLayoutCftAfterLoFTitle}%
  \setlength{\cftafterlottitleskip}{\ITBLayoutCftAfterLoTTitle}%
}

% ---------------------------------------------------------------------------
% (4) Sampul & halaman pengesahan — vertical gaps (juknis diagram)
% ---------------------------------------------------------------------------
\newcommand{\ITBSampulTopGap}{0.2cm}
\newcommand{\ITBSampulAfterTitle}{1.0cm}
\newcommand{\ITBSampulAfterDisertasi}{0.5cm}
\newcommand{\ITBSampulAfterKaryaTulis}{0.9cm}
\newcommand{\ITBSampulAfterOleh}{0.3cm}
\newcommand{\ITBSampulAfterNimBlock}{0.2cm}
\newcommand{\ITBSampulAfterProdi}{0.9cm}
\newcommand{\ITBSampulLogoHeight}{3.5cm}
\newcommand{\ITBSampulAfterInstitusi}{0.15cm}

\newcommand{\ITBPengesahanAfterHeading}{0.8cm}
\newcommand{\ITBPengesahanAfterTitle}{0.8cm}
\newcommand{\ITBPengesahanAfterAuthorBlock}{0.9cm}
\newcommand{\ITBPengesahanAfterMenyetujui}{0.3cm}
\newcommand{\ITBPengesahanAfterTanggal}{0.6cm}
\newcommand{\ITBPengesahanAfterKetuaLabel}{1.5cm}
\newcommand{\ITBPengesahanAfterKetuaRule}{0.8cm}
\newcommand{\ITBPengesahanAnggotaRuleDrop}{1.5cm}

\newcommand{\ITBPedomanBetweenParagraphs}{0.5cm}
|.
%% Compile from the directory that contains this file (same folder as
%% |disertasi.tex|) so the input resolves.
%%
%% Sections:
%%   (1) Page geometry — mirror margins (Pedoman / Juknis SPs ITB)
%%   (2) Heading vertical rhythm — titlesec (chapter aligns to 3 cm rule)
%%   (3) TOC / LOF / LOT — tocloft skips and list title vertical tweak
%%   (4) Sampul & halaman pengesahan — vertical gaps (juknis diagram)
%%
%% To adjust globally: change the macro or skip value here only, then
%% rebuild. For empirical fine-tuning of chapter vs.\ topskip offset,
%% edit |\ITBLayoutChapterTitleAbove| (negative = pull title up).
%% ---------------------------------------------------------------------------

% ---------------------------------------------------------------------------
% (1) Page geometry — applied via |\ITBApplyPageGeometry| after
%     |\RequirePackage{geometry}| in the class file.
% ---------------------------------------------------------------------------
\newcommand{\ITBGeomPaper}{a4paper}
\newcommand{\ITBGeomInner}{4cm}
\newcommand{\ITBGeomOuter}{3cm}
\newcommand{\ITBGeomTop}{3cm}
\newcommand{\ITBGeomBottom}{3cm}
\newcommand{\ITBGeomHeadHeight}{15pt}
\newcommand{\ITBGeomHeadSep}{12pt}
\newcommand{\ITBGeomFootSkip}{30pt}

\newcommand{\ITBApplyPageGeometry}{%
  \geometry{%
    \ITBGeomPaper,%
    inner=\ITBGeomInner,%
    outer=\ITBGeomOuter,%
    top=\ITBGeomTop,%
    bottom=\ITBGeomBottom,%
    headheight=\ITBGeomHeadHeight,%
    headsep=\ITBGeomHeadSep,%
    footskip=\ITBGeomFootSkip,%
    twoside%
  }%
}

% ---------------------------------------------------------------------------
% (2) Chapter / section vertical spacing (titlesec + report baseline).
%     Chapter: single line "Bab I  Judul" per juknis Mei 2019 hal.\ 6.
% ---------------------------------------------------------------------------
\newskip\ITBLayoutChapterTitleAbove
\newskip\ITBLayoutChapterTitleBelow
\newskip\ITBLayoutSectionTitleAbove
\newskip\ITBLayoutSectionTitleBelow
\newskip\ITBLayoutSubsectionTitleAbove
\newskip\ITBLayoutSubsectionTitleBelow
\newskip\ITBLayoutSubsubsectionTitleAbove
\newskip\ITBLayoutSubsubsectionTitleBelow

\ITBLayoutChapterTitleAbove=-36pt plus 0pt minus 0pt
\ITBLayoutChapterTitleBelow=18pt plus 0pt minus 0pt
\ITBLayoutSectionTitleAbove=18pt plus 0pt minus 0pt
\ITBLayoutSectionTitleBelow=6pt plus 0pt minus 0pt
\ITBLayoutSubsectionTitleAbove=14pt plus 0pt minus 0pt
\ITBLayoutSubsectionTitleBelow=4pt plus 0pt minus 0pt
\ITBLayoutSubsubsectionTitleAbove=12pt plus 0pt minus 0pt
\ITBLayoutSubsubsectionTitleBelow=4pt plus 0pt minus 0pt

% ---------------------------------------------------------------------------
% (3) Daftar isi / gambar / tabel — tocloft vertical density + title skip
% ---------------------------------------------------------------------------
\newcommand{\ITBLayoutCftBeforeChapSkip}{0.2em}
\newcommand{\ITBLayoutCftBeforeSecSkip}{0.1em}
\newcommand{\ITBLayoutCftBeforeSubsecSkip}{0.1em}
\newcommand{\ITBLayoutCftBeforeFigSkip}{0.1em}
\newcommand{\ITBLayoutCftBeforeTabSkip}{0.1em}
\newlength{\ITBLayoutCftBeforeTocTitle}
\newlength{\ITBLayoutCftBeforeLoFTitle}
\newlength{\ITBLayoutCftBeforeLoTTitle}
\newlength{\ITBLayoutCftAfterTocTitle}
\newlength{\ITBLayoutCftAfterLoFTitle}
\newlength{\ITBLayoutCftAfterLoTTitle}

\setlength{\ITBLayoutCftBeforeTocTitle}{-20pt}
\setlength{\ITBLayoutCftBeforeLoFTitle}{-20pt}
\setlength{\ITBLayoutCftBeforeLoTTitle}{-20pt}
\setlength{\ITBLayoutCftAfterTocTitle}{18pt}
\setlength{\ITBLayoutCftAfterLoFTitle}{18pt}
\setlength{\ITBLayoutCftAfterLoTTitle}{18pt}

\newcommand{\ITBApplyTocloftLayout}{%
  \setlength{\cftbeforechapskip}{\ITBLayoutCftBeforeChapSkip}%
  \setlength{\cftbeforesecskip}{\ITBLayoutCftBeforeSecSkip}%
  \setlength{\cftbeforesubsecskip}{\ITBLayoutCftBeforeSubsecSkip}%
  \setlength{\cftbeforefigskip}{\ITBLayoutCftBeforeFigSkip}%
  \setlength{\cftbeforetabskip}{\ITBLayoutCftBeforeTabSkip}%
  \setlength{\cftbeforetoctitleskip}{\ITBLayoutCftBeforeTocTitle}%
  \setlength{\cftbeforeloftitleskip}{\ITBLayoutCftBeforeLoFTitle}%
  \setlength{\cftbeforelottitleskip}{\ITBLayoutCftBeforeLoTTitle}%
  \setlength{\cftaftertoctitleskip}{\ITBLayoutCftAfterTocTitle}%
  \setlength{\cftafterloftitleskip}{\ITBLayoutCftAfterLoFTitle}%
  \setlength{\cftafterlottitleskip}{\ITBLayoutCftAfterLoTTitle}%
}

% ---------------------------------------------------------------------------
% (4) Sampul & halaman pengesahan — vertical gaps (juknis diagram)
% ---------------------------------------------------------------------------
\newcommand{\ITBSampulTopGap}{0.2cm}
\newcommand{\ITBSampulAfterTitle}{1.0cm}
\newcommand{\ITBSampulAfterDisertasi}{0.5cm}
\newcommand{\ITBSampulAfterKaryaTulis}{0.9cm}
\newcommand{\ITBSampulAfterOleh}{0.3cm}
\newcommand{\ITBSampulAfterNimBlock}{0.2cm}
\newcommand{\ITBSampulAfterProdi}{0.9cm}
\newcommand{\ITBSampulLogoHeight}{3.5cm}
\newcommand{\ITBSampulAfterInstitusi}{0.15cm}

\newcommand{\ITBPengesahanAfterHeading}{0.8cm}
\newcommand{\ITBPengesahanAfterTitle}{0.8cm}
\newcommand{\ITBPengesahanAfterAuthorBlock}{0.9cm}
\newcommand{\ITBPengesahanAfterMenyetujui}{0.3cm}
\newcommand{\ITBPengesahanAfterTanggal}{0.6cm}
\newcommand{\ITBPengesahanAfterKetuaLabel}{1.5cm}
\newcommand{\ITBPengesahanAfterKetuaRule}{0.8cm}
\newcommand{\ITBPengesahanAnggotaRuleDrop}{1.5cm}

\newcommand{\ITBPedomanBetweenParagraphs}{0.5cm}
